\documentclass[10pt,a4paper]{amsart}
\usepackage[margin=19mm]{geometry}
\usepackage{amsmath,amssymb,mathtools}
\usepackage[scr=rsfs,cal=euler]{mathalfa}
\usepackage{xfrac}
\usepackage[unicode,hidelinks,bookmarks=false]{hyperref}
\newcommand{\ie}{i.e.\ }
\newcommand{\cf}{cf.\ }
\newcommand{\resp}{respectively\ }
\newcommand{\Subsection}{\subsection}
\providecommand{\nobreakdash}{\nobreak\hbox{-}}
\setlength{\emergencystretch}{2em}
\multlinegap0pt
\usepackage{bm}
\usepackage{bbm}
\usepackage{dsfont}
\usepackage{xspace}
\usepackage{cancel}
\usepackage{mathtools}
\usepackage{amsthm}
\makeatletter
\@ifundefined{thm}{\newtheorem{thm}{Thm}}{}
\@ifundefined{Fact}{\newtheorem{Fact}{Fact}}{}
\@ifundefined{corollary}{\newtheorem{corollary}[thm]{Corollary}}{}
\@ifundefined{definition}{\newtheorem{definition}[thm]{Definition}}{}
\@ifundefined{lemma}{\newtheorem{lemma}[thm]{Lemma}}{}
\makeatother
\def\dom{\operatorname{dom}}
\def\val{\operatorname{val}}
\title[Not OCA and products of Frechet spaces]{Not OCA and products of Frechet spaces}
\author{Alan Dow}
\date{Source: arXiv:2507.01662v1 (CC BY 4.0)}
\begin{document}
\maketitle

\noindent\emph{Source and licence.} Not OCA and products of Frechet spaces. Version arXiv:2507.01662v1 (CC BY 4.0), distributed under the Creative Commons Attribution 4.0 International licence, as recorded in the arXiv licence metadata for this exact version and shown on the abstract page https://arxiv.org/abs/2507.01662v1. Full rights notice: licenses/arxiv-2507.01662-CC-BY-4.0.txt.\par\smallskip

\noindent\textit{Editorial scope.} Counted unit: the Lemma whose source label is \texttt{cohere} in the article (source0.tex lines 1139--1167, source label \texttt{cohere}), delivered together with the article's own complete original proof of it (source0.tex lines 1169--1419, one \texttt{proof} environment of 251 lines). No mathematics is written, repaired or re-derived: statement and proof are the article's own text line for line. Required context retained uncounted: defcoherent (lines 948--972); successor (lines 1016--1019); cccequivalent (lines 1042--1048); Q' (lines 1102--1137). Every definition, display and label of those lines is retained unchanged. Omitted from this package and not redistributed: 1--947, 973--1015, 1020--1041, 1049--1101, 1138--1138, 1168--1168, 1420--1827. Nothing inside a retained excerpt is omitted or altered: every retained body is delivered exactly as the article's lines are written, so no omission receipt of any kind accompanies this package. Citations: the retained excerpts carry no citation command at all, so no bibliography entry of the article is transcribed and no reference is redistributed. Reference adaptation: none was needed and none was made; every reference inside the delivered unit resolves inside it, so no unresolved reference or citation is produced. Printed slips kept literal: no printed word, symbol, hypothesis or number of the counted statement or of its proof was altered. Layout and class: the article's own class is \texttt{amsart} with packages including amssymb, amsrefs; the wrapper is the lane's pinned \texttt{amsart} editorial wrapper at 10pt on A4, which restates the theorem-like environments the retained text needs and adds the article's own macro definitions that the retained text needs (\texttt{\textbackslash dom}, \texttt{\textbackslash val}), with extra wrapper packages (bm, bbm, dsfont, xspace, cancel, mathtools). The delivered numbering is the unit's own. Assets: no retained span contains a figure environment, a graphics-inclusion command, a PDF-page inclusion, a table or a TikZ picture; no image file is copied into this package, the package declares no asset dependency and no image file is redistributed with it. Licence: version arXiv:2507.01662v1 of this article is distributed under the Creative Commons Attribution 4.0 International licence, as recorded in the arXiv licence metadata for this exact version and shown on the abstract page https://arxiv.org/abs/2507.01662v1. A full rights notice is delivered at licenses/arxiv-2507.01662-CC-BY-4.0.txt. Characters: every retained excerpt is pure ASCII, so the delivered engine, XeLaTeX with its default Latin Modern text font, reports no missing character.
\begin{definition}  If $\mathcal H_\lambda$ is a linear\label{defcoherent} coherent family, 
then $Q(\mathcal H_\lambda)$, also denoted   $Q(\{ h_\alpha : \alpha < \lambda\})$,
is the following poset. A condition $q\in Q(\mathcal H_\lambda)$, is
 a tuple $(s_q, h_q, F_q, f_q)$ satisfying 
 \begin{enumerate}
\item   $s_q\in \omega^{<\omega}$ with domain $n_q$ and $f_q\in \omega^\omega$,
\item $h_q $  is a 2-valued  function with domain  $s_q^\downarrow $,
 \item   $F_q$ is a finite subset of $\lambda$,
 \item for all $\alpha\in F_q$ and $\delta_q = \max(F_q)$,
 and for all $n_q\leq m$,
  $f_\alpha(m)\leq f_{\delta_q}(m) \leq f_q(m)$,
  and for all $(m,j)\in \dom(h_\alpha)\cap \dom(h_{\delta_q})$,
  $h_\alpha(m,j) = h_{\delta_q}(m,j)$.
 \end{enumerate}
 The ordering on $Q(\mathcal H_\lambda)$ is that $q\leq r$
 providing $s_q\supset r_q$, $h_q\supset h_r$, $F_q\supset F_r$,
  $f_q\geq f_r$, and for all $(m,j)\in \dom(h_q)\cap \dom(h_{\delta_r})$
  with $n_r\leq m$, $s_q(m,j) = h_{\delta_r}(m,j)$. 

  We let $\dot f$ and $\dot h$ be the two canonical $Q(\mathcal H_\lambda)$-names,
    (and in context we would  denote them as $\dot f_\lambda$ and $\dot h_\lambda$)
    where, if $G$ is a $Q(\mathcal H_\lambda)$-generic filter,
     $\val_{G}(\dot f) = \bigcup \{ s_q : q\in G\}$
     and $\val_{G}(\dot h) = \bigcup\{ h_q : q\in G\}$. 
\end{definition}

\begin{corollary}  If $\{ h_\alpha : \alpha \leq  \lambda\}$ 
is a linear coherent gap, then\label{successor} 
 $Q(\{ h_\alpha : \alpha < \lambda\})$ is $\sigma$-centered. 
\end{corollary}

  \begin{lemma} For  an uncountable\label{cccequivalent}
  linear coherent gap, $\{ h_\alpha : \alpha<\lambda\}$,
the poset   $  Q(\{ h_\alpha :\alpha < \lambda\})$ is ccc 
if and only if 
for every uncountable $X\subset\lambda$, there are
$\alpha< \beta$ in $X$ such that $h_\alpha\cup h_\beta$ is a function.
  \end{lemma}

\begin{definition} For a poset $P$ and   $P$-names, 
 $\{ \dot f_\alpha, \dot h_\alpha : \alpha < \lambda\}$,
 that is forced to be a linear coherent sequence, we define\label{Q'} 
 the $P$-name $\dot Q'(\{ \dot h_\alpha : \alpha < \lambda\})$ as
 follows. A condition $q\in 
 \dot Q'(\{ \dot h_\alpha : \alpha < \lambda\})$ is a tuple
  $( \check{n}_q, \check{s}_q, \tau_q, \check{h}_q, \pi_q, \check{F_q}, \dot f_q)$
  where the following are forced by $1_{P}$: 
  \begin{enumerate}
  \item  $n_q\in \omega$, $s_q   \in \omega^{n_q}$, $s_q\leq \tau_q\in \omega^{n_q}$,
  \item $h_q$ is a 2-valued function with domain $s_q^{\downarrow}$,
  \item $h_q\subset \pi_q$ is a 2-valued function with domain $\tau_q^{\downarrow}$,
  \item $F_q$ is a finite subset of $\lambda$,
  \item for each $\alpha\in F_q$, $\dot f_\alpha \leq \dot f$.
  \end{enumerate}
Say that a condition $q\in \dot Q'(\{ \dot h_\alpha : \alpha < \lambda\})$ 
is pure if $\tau_q = \check{s}_q $ and $\pi_q = \check{h}_q$.



For each $q\in \dot Q'(\{\dot h_\alpha :\alpha < \lambda\})$, 
  let $\hat q = ( \tau_q , \pi_q, F_q, \dot f_q)$.  We note that
   $\hat q$ is forced to be an element of $Q(\{ \dot h_\alpha :\alpha < \lambda\})$
   and we define the ordering on 
$    \dot Q'(\{\dot h_\alpha :\alpha < \lambda\})$ 
by $q_1 \leq q_2$ if 
  $\hat{q}_1 \leq \hat{q}_2 $.
\medskip

 Suppose that $s\in \omega^n$ and $\tau$ is a $P$-name for an
 element of $\omega^n$ such that $1\Vdash s\leq \tau$. 
 For any 2-valued function $h$ with domain $s^{\downarrow}$,
  let $h\oplus 0_{\tau}$ denote the name of the function with domain
   $\tau^{\downarrow}$ that extends $h$ and has value $0$
   at all $(m,j)\in \tau^{\downarrow}\setminus s^\downarrow$.
\end{definition}

\begin{lemma}  Suppose that $\langle P_\alpha, \dot Q_\beta : \alpha \leq \lambda,
\beta <\lambda\rangle$ is a finitely\label{cohere} supported iteration,
and $\{ \dot f_\alpha , \dot h_\alpha : \alpha < \lambda \}$ are
 $P_\lambda$-names 
satisfying
\begin{enumerate}
\item for each $\alpha<\lambda$, each of $\dot f_\alpha$ and 
 $\dot h_\alpha$ are $P_{\alpha+1}$-names  satisfying that $\dot f_\alpha$
 is forced to be in $\omega^\omega$ and $\dot h_\alpha$ is a 2-valued
 function with domain $\dot f_\alpha^\downarrow$,
\item for each $\beta<\lambda$, $\dot Q_\beta$ is a $P_\beta$-name of a ccc poset,
\item for each even ordinal $\beta<\lambda$, if $P_\beta$ forces
that $\{ \dot h_\alpha : \alpha < \beta \}$ 
is a linear coherent family, then $\dot Q_\beta$ is 
the $P_\beta$-name  $\dot Q'(\{\dot h_\alpha : \alpha <\beta \})$,
and $\dot f_\beta, \dot h_\beta$ are the canonical $P_{\beta+1}$-names
 associated with $\dot Q_\beta$, otherwise $\dot Q_\beta = 2^{<\omega}$ and
 $\dot f_\beta=\dot f_0$,
 and $\dot h_\beta = \dot h_0$,
 \item if $\beta=\alpha+1<\lambda$ is an odd ordinal, then $\dot f_\beta=\dot f_\alpha$
 and $\dot h_\beta = \dot h_\alpha$.
 \end{enumerate}
Then for each $\beta \leq \lambda$, $P_\beta$ forces
that $\{ \dot h_\alpha : \alpha < \beta \}$ is
a linear coherent family. Furthermore, if $\beta$ has uncountable cofinality,
 then $P_\beta$ forces that for every uncountable $X\subset \beta$,
  there are distinct $\xi,\eta\in X$ such that
   $\dot h_\xi\cup \dot h_\eta$ is a function.
\end{lemma}

  \begin{proof}  It is immediate from the remarks immediately 
  after Definition \ref{defcoherent} that, for each $\beta <\lambda$,
   $P_\beta$ forces that $\{ \dot h_\alpha : \alpha < \beta\}$ is
    a linear coherent family. It, however, is not immediate
    that $P_\beta$ is ccc. We prove the second stated conclusion   
    of the Lemma by induction on $\beta \leq \lambda$. 
Since this conclusion is vacuous for $\beta < \omega_1$, we
might as well simply assume that it holds for all $\beta < \lambda$
and prove it for $\lambda$.  
It follows from Corollaries \ref{successor} and \ref{cccequivalent}
that for uncountable $X\subset \beta <\lambda$, it remains
true, i.e. not destroyed by further forcing,
that there are $\xi<\eta\in X$ such that
 $h_\xi\cup h_\eta$ is a function.
Needless to say, there is nothing to 
prove unless $\lambda$ has cofinality $\omega_1$. 

In preparation we make some observations, stated as Facts, about $P_\lambda$.

\begin{Fact} Let $ \tilde P_\lambda$ be the set of conditions that satisfy, for each
even $\beta\in \dom(p)$ and each odd $\alpha+1\in F_{p(\beta)}$,
 we also have that $\alpha$ is in $F_{p(\beta)}$.  Then
  $\tilde P_\lambda$ is a dense subset of $P_\lambda$.
\end{Fact}

 

\begin{Fact} Let   $p\in \tilde P_\lambda$ and suppose\label{fullfour}
that $\alpha <\beta$ are even ordinals in $\dom(p)$. Then $p$ will force
that $\dot h_\alpha\cup \dot h_\beta$ is a function if
\begin{enumerate}
\item $\alpha \in F_{p(\beta)}$,
\item $n_{p(\alpha)} = n_{p(\beta)}$, $s_{p(\alpha)} = s_{p(\beta)}$,
   $h_{p(\alpha)} = h_{p(\beta)}$,
   \item $\pi_{p(\alpha)} = h_{p(\alpha)} \oplus 0_{\tau_{p(\alpha)}}$ (as in Definition \ref{Q'}), and
   \item $\pi_{p(\beta)} = h_{p(\beta)} \oplus 0_{\tau_{p(\beta)}}$.
\end{enumerate}
\end{Fact}


\begin{Fact} Let   $p\in \tilde P_\lambda$ and suppose\label{addF}
that $\alpha <\beta$ are even ordinals in $\dom(p)$ and
that $\alpha\in F_{p(\beta)}$. Then $\bar p$ is an extension of $p$
where $\bar p(\gamma) = p(\gamma)$ for all $\beta\neq\gamma\in \dom(p)$
and 
 $\bar p(\beta)$ is equal to
 $
 ( n_{p(\beta)}, s_{p(\beta)}, \tau_{p(\beta)}, 
  h_{p(\beta)}, 
     \pi_{p(\beta)},   F_{p(\beta)}\cup F_{p(\alpha)}, \dot f_{p(\beta)})$.
\end{Fact}

\begin{Fact} Suppose that $p\in \tilde P_\lambda$ and that for
even $\alpha < \beta$ both in $\dom(p)$ we have\label{just3}
the conditions (2)-(4) of Fact \ref{fullfour} holding
and that $\delta_{p(\alpha)}$ is an element of $F_{p_\xi(\beta)}$.
Then the  condition $\bar p \leq p$ forces that $
\dot h_\alpha\cup \dot h_\beta$ is a function where
$\bar p$ is defined as follows: 
$\bar p(\gamma)=p(\gamma)$ for all $\beta\neq \gamma \in \dom(p)$,
and 
$$\bar p(\beta) = (n_{p(\beta)}, s_{p(\beta)},
 \tau_{p(\beta)}, h_{p(\beta)}, h_{p(\beta)}\oplus 0_{\tau_{p(\beta)}} , 
   F_{p(\beta)}\cup \{ \alpha\}, \dot f_{p(\beta)} \vee \dot f_\alpha)~.$$
\end{Fact}

\begin{Fact} Let $p\in \tilde P_\lambda$ and\label{addmu}
let $  \beta$ be an even ordinal
in $\dom(p)$ such that  $p(\beta)$ is a   pure condition
and let $\delta = \delta_{p(\beta)}$ be the maximum even ordinal 
in $  \max(F_{p(\beta)})$.  Choose
any $m\in\omega$ such that $n = n_{p(\beta)}< m$ and any
$\bar p\in P_{\delta+1}$ that forces   values $\bar s, \bar h$
on $\dot f_\delta\restriction [n,m)$ and $\dot h_\delta\restriction 
\bar s^{\downarrow}$ respectively.
Recall that $1_{P_\beta}$ forces that $\dot f_\delta\leq \dot f_{p(\beta)}$.
Then $\tilde p$ is an extension
of $\bar p$ and $p$ where $\tilde p\restriction \delta = \bar p$,
 $\tilde p(\gamma) = p(\gamma)$ for $\beta\neq \gamma \in \dom(p) 
 \setminus \delta{+}1$,
 and $\tilde p(\beta) =
 (m, s_{p(\beta)}\cup \bar s, \bar \tau , h_{p(\beta)}\cup \bar h,
 ( h_{p(\beta)}\cup \bar h)\oplus 0_{\bar\tau}, F_{p(\beta)}, \dot f_{p(\beta)})$
 where $\bar\tau = s_{p(\beta)}\cup \dot f_{p(\beta)}\restriction [n,m)$.
 \medskip

 Moreover, if for some even $\alpha <\beta$, $\bar p$ forces
 that $\dot h_\delta\cup (\dot h_\alpha \restriction \left([m,\omega)\times \omega )\right)$
 is a function, then we could instead define $\tilde p(\beta)$
 to equal $
 (m, s_{p(\beta)}\cup \bar s, \bar \tau , h_{p(\beta)}\cup \bar h,
 ( h_{p(\beta)}\cup \bar h)\oplus 0_{\bar\tau}, F_{p(\beta)}\cup\{\alpha\}, \dot f_{p(\beta)}
\vee \dot f_\alpha )$
\end{Fact}

 
\bigskip

With the benefit of the above Facts we are ready to prove Theorem \ref{cohere}.
Let $\dot X$
be a $P_\lambda$-name of an uncountable subset of $\lambda$.
By the definition of the family $\{ \dot h_\alpha : \alpha<\lambda\}$,
 we may assume that $\dot X$ is forced to consist of even ordinals.
Since we are proceeding by induction, we may assume
that $\dot X$ is forced to be cofinal in $\lambda$.
Let $e$ be a strictly increasing function from $\omega_1$
to a cofinal subset of $\lambda$.
 For each $\xi < \omega_1$, choose a condition $p_\xi\in \tilde P_\lambda$
 that forces some $ \beta_\xi \in \lambda\setminus e(\xi) $ is an element
 of $\dot X$. For each $\xi\in\omega_1$, let $\beta_\xi\in H_\xi$ denote the finite
 support of $p_\xi$. 

 For each $\xi\in\omega_1$, we make some additional assumptions
 about $p_\xi$. Let $\beta$ be the maximum even ordinal in $H_\xi$. 
 By possibly strengthening $p_\xi\restriction \beta$ we can ensure
 that $p_\xi\restriction \beta$ forces that $p_\xi(\beta)$ is pure. 
 We can also ensure that $F_{p_\xi(\beta)} \cap \alpha $ is a subset
 of $\dom(p_\xi)$, and  for each even $\alpha\in \dom(p_\xi)\cap \beta$,
 $F_{p_\xi(\beta )}\cap \alpha $ is a subset of $F_{p_\xi(\alpha)}$.   
 We can also ensure that $n_{p_\xi(\alpha)}\geq  n_{p_\xi(\beta)}  $ 
 for all even $\alpha\in \dom(p_\xi)$.
 This is step 1 of a finite recursion (since every descending
 sequence of ordinal is finite).  In this way we can 
 assume that, for each even ordinal $\beta$ in $\dom(p_\xi)$,
  $p_\xi\restriction \beta$ forces that $p_\xi(\beta)$ 
  is pure 
 and for even $\alpha\in H_\xi\cap \beta$,
 $F_{p_\xi(\beta)}\cap \alpha \subset
  F_{p_\xi(\alpha)}$.
  
 
 By passing to an uncountable subset we can 
 assume that each $H_\xi$ has cardinality $\ell$ and fix
 an increasing enumeration, $\{ \alpha(\xi,i) : i<\ell\}$ of $H_\xi$. 
 For each $i<\ell$ and $\xi,\eta<\omega_1$, we may assume
 that $\alpha(\xi,i)$ is even if and only if $\alpha(\eta,i)$ is even,
 and that  there is a fixed $\bar \imath<\ell$ so
 that $\beta_\xi  = \alpha(\xi, \bar \imath)$ for all $\xi$. 
 Let $E$ denote the set of $i<\ell$ such that (each)
  $\alpha(\xi,i)$ is even
We can also assume that for $\xi <\eta$ and $i\in E $,
$(n_i, s_i, h_i) =  
 (n_{p_\xi( \alpha(\xi,i))} , s_{p_\xi( \alpha(\xi,i))} , h_{p_\xi( \alpha(\xi,i))} ) = 
(n_{p_\xi( \alpha(\xi,i))} , s_{p_\eta(\alpha(\eta,   i))} , h_{p_\eta(\alpha(\eta,i))})
$.
 
 Notice that $\{ \alpha(\xi,\bar \imath) : \xi\in\omega_1\}$
 is unbounded in $\lambda$. 
 By a recursion of length at most $\bar \imath$, we can repeatedly
 pass to an uncountable subset of $\xi\in\omega_1$ so as to ensure,
  for each $i<\bar \imath$, either $\{ \alpha(\xi,i) : \xi\in\omega_1\}$
   is unbounded in $\lambda$ or has an upper bound $\mu_i<\lambda$.
   Let $i_0 \leq  \bar \imath$ be minimal so that 
    $\{ \alpha(\xi,i_0) : \xi\in\omega_1\}$ is unbounded. 
    Choose any even $\mu<\lambda$ so that $\alpha(\xi,i)<\mu$
    for all $\xi\in\omega_1$ and $i<i_0$. 
    For simple convenience assume that,
    for all $\xi\in\omega_1$, $\alpha(\xi,i_0)$ is an even ordinal. 
 Let   $\{ i_k : k <\bar\ell\} $ be an enumeration of
  $E\setminus i_0$.

By the inductive assumption, $P_{\mu+1}$ is ccc and so we may
choose a generic filter $G_{\mu+1}$ for $P_{\mu+1}$ such
that $\Gamma = \{ \xi \in\omega_1 : p_\xi \restriction \mu \in G_\mu\}$ is
uncountable. By re-indexing we may assume that $\mu+1< \alpha(\xi,i_0)$
for all $\xi\in\omega_1$, and by again choosing an
uncountable subsequence
 we can assume that, for $\xi<\eta$ both in $\Gamma$,
   $H_\xi \subset \alpha(\eta,i_0)$.

For each $\xi\in\Gamma$, let $\delta_\xi$ denote the maximum
element of $F_{p_\xi(\alpha(\xi,i_0))}$.  Since $\delta_\xi<\mu$
we may  
let $f_{\delta_\xi}$ and $h_{\delta_\xi}$ be the valuations
of $\dot f_{\delta_\xi}$ and  $\dot h_{\delta_\xi}$ respectively
by the filter $G_{\mu+1}$. Let
 $f_\mu$ and $h_\mu$ be defined analogously.
Now choose a value $\bar m\in\omega$ and an uncountable
 $\Gamma_1\subset \Gamma$ satisfying that
  $f_{\delta_\xi} \restriction [\bar m,\omega)\leq f_\mu$,
  and 
   $h_{\delta_\xi}\restriction [\bar m,\omega)\times\omega$
   is a subset of $h_\mu$ for all $\xi\in \Gamma_1$.

We work in the extension $V[G_{\mu+1}]$ and fix any $\xi\in \Gamma_1$.
We can ignore $p_\xi\restriction \mu$ since we have
that $p_\xi\restriction \mu\in G_{\mu+1}$.  
 For each even $\beta\in H_\xi\setminus \mu$, 
 we have that  $\delta_\xi$ is an
 element of $F_{p_\xi(\beta)}$ and is
 the maximum of $F_{p_\xi(\beta)}\cap \mu+1$.
 Let $\bar s = f_{\delta_\xi}\restriction [n_{i_0} , \bar m)$
 and $\bar h = h_{\delta_\xi}\restriction \bar s^\downarrow$. 
Applying the ``moreover'' clause of Fact \ref{addmu}, we have an extension
 $\bar p_\xi$ of $p_\xi$ such
 that $\bar p_\xi \restriction \mu{+}1\in G_{\mu+1}$ forces
 that
 $\bar s = \dot f_{\delta_\xi}\restriction [n_{i_0} , \bar m)$
 and $\bar h = \dot h_{\delta_\xi}\restriction \bar s^\downarrow$,
 $\bar p_\xi(\gamma) = p_\xi(\gamma)$ for $\alpha(\xi,i_0)< \gamma \in H_\xi$
 and
  $\bar p_\xi(\alpha(\xi,i_0))$ is as indicated in Fact \ref{addmu}
  with $\mu$ added to $ F_{\bar p_\xi(\alpha(\xi,i_0))}$.
  Next, by a finite recursion, we keep applying the first
  clause of Fact \ref{addmu}, so as to arrange that
    $\alpha(\xi,i_k)\in F_{\bar p_\xi(\alpha(\xi,i_{k+1}))}$
    and
     $n_{\bar p_\xi(\alpha(\xi,i_k))}=n_{\bar p_\xi(\alpha(\xi,i_{k+1}))}$
     for all
     $  k<\bar \ell-1$.  
     Actually we can stop at
      $\alpha(\xi,\bar\imath ) $ but nothing is saved.
      Additionally, by Fact \ref{addF}, we can assume
      that $F_{\bar p_\xi(\alpha(\xi,i_k))} \subset
         F_{\bar p_\xi(\alpha(\xi,i_{k+1}))}$ for all $k<\bar\ell$.

 Choose $\xi <\eta$ both in $\Gamma_1$
 so that, for all $k < \bar \ell-1$,
  $$( n_{\bar p_\xi(\alpha(\xi,i_k))} , 
 s_{\bar p_\xi(\alpha(\xi,i_k))} , 
  h_{\bar p_\xi(\alpha(\xi,i_k))})
  =( n_{\bar p_\eta(\alpha(\xi,i_k))} , 
 s_{\bar p_\eta(\alpha(\xi,i_k))} , 
  h_{\bar p_\eta(\alpha(\xi,i_k))})
  ~.$$
 Since $\bar p_\xi \restriction \mu+1$ and
 $\bar p_\eta\restriction \mu+1$ are both in $G_{\mu+1}$,
 and $H_\xi\cap H_\eta\subset \mu$, there is a condition
  $\bar p\in P_\lambda$ satisfying that
   $\bar p\setminus \mu+1 = (\bar p_\xi \restriction H_\xi\setminus(\mu+1))
   \cup (\bar p_\eta \restriction H_\eta \setminus (\mu+1))$.
 
  Recall that $\beta_\xi = \alpha(\xi, \bar\imath)$.
Note that $\mu$ is the largest element of $F_{\bar p_\eta(\alpha(\eta,i_0))}$
and that $\bar p_\xi$ forces that
     $ h_\mu\restriction [n_{\bar p_\xi(\beta_\xi)},\omega)
     \times\omega)$ is a subset of $\dot h_{\beta_\xi}$.     
     Therefore, by the \textit{moreover} clause of Fact \ref{addmu},
     we can extend $\bar p $ 
     to a condition $p'$ (but only in the coordinate
      $\alpha(\eta,i_0)$) as in Fact \ref{addmu} so
      that $F_{p'(\alpha(\eta,i_0))}$ is obtained
      by   adding  $\beta_\xi$ to $F_{\bar p_\eta(\alpha(\eta,i_0))}$,
      and, by Fact \ref{addF}, we can also arrange
      that $\beta_\xi $ is an element
      of $F_{  p' (\beta_\eta)}$. 
      The proof is finished by verifying
      that $\alpha  = \beta_\xi $ and $\beta = \beta_\eta$
      satisfy the conditions of Fact \ref{fullfour} for
      the condition $p'$.
         \end{proof}

\end{document}
